\section{The Wage Paid, the Enemy Still in the Field}

The New Testament speaks about death in two tenses, and the awkwardness is on the surface of the text rather than hidden in it.

In one tense the thing is finished. Christ ``hath destroyed death and hath brought to light life and incorruption by the gospel.''\endnote{2 Timothy 1:10, Douay--Rheims, read 2026-07-26 in the tracked Challoner verse text. The verse's participle governs a completed action: the destruction is spoken of as accomplished and made manifest.} He became a partaker of flesh and blood ``that, through death, he might destroy him who had the empire of death, that is to say, the devil: and might deliver them, who through the fear of death were all their lifetime subject to servitude.''\endnote{Hebrews 2:14--15, Douay--Rheims, read as above. The word rendered ``servitude'' is worth pausing on: the writer's diagnosis is not that men are afraid at the end but that the fear operates over a whole lifetime and produces a condition of bondage---the fear, not the event, is what enslaves. This essay reads the verse for that diagnosis and does not develop the atonement theology in which it stands.}

In the other tense the thing is pending. Paul, laying out the order of the resurrection, arrives at the end of the sequence and names what is left: ``For he must reign, until he hath put all his enemies under his feet. And the enemy, death, shall be destroyed last.''\endnote{1 Corinthians 15:25--26, Douay--Rheims, read as above. The Douay's word order---``the enemy, death, shall be destroyed last''---keeps the appositive; the more familiar English is ``the last enemy that shall be destroyed is death.'' The immediate context is Paul's ordering at 15:23--24: ``But every one in his own order: the firstfruits, Christ: then they that are of Christ, who have believed in his coming. Afterwards the end.'' The verse's placement matters for the argument here: death is named at the end of a queue, which presupposes that it is still standing.} Last. Which is to say: still standing, still an enemy, and still to be dealt with at a date not yet arrived.

Both statements are in the same canon and neither is a slip. The difficulty they create is not a difficulty about words. It is the ordinary experience of every Christian who has ever buried one, and the earliest generations felt it more sharply than we do, because they were being asked to explain it to people watching them die of the same plagues as everyone else.

\subsection{The complaint from the plague}

Cyprian of Carthage wrote to his people during an epidemic, and the objection he answers is not a theological subtlety. It is the obvious one. ``But nevertheless it disturbs some that the power of this Disease attacks our people equally with the heathens, as if the Christian believed for this purpose, that he might have the enjoyment of the world and this life free from the contact of ills.''\endnote{Cyprian of Carthage, \emph{De mortalitate} 8, in the English of the Ante-Nicene Fathers, vol.~V (Buffalo: Christian Literature Publishing Company, 1886), tr.\ Robert Ernest Wallis; read 2026-07-26 in the New Advent transcription of that volume (Treatise 7). The treatise was occasioned by the epidemic that struck Carthage from 252; this essay uses it as a patristic and pastoral witness, not as a doctrinal authority, and did not collate the English against Cyprian's Latin.}

His answer refuses every available comfort. He does not say the Christians are being spared, because they are not. He does not say the disease is a special providence for the unbelievers, because it is not. He says that the flesh is common: ``So long as we are here in the world, we are associated with the human race in fleshly equality, but are separated in spirit.'' Then a list which sounds, at a distance of seventeen centuries, like a man refusing to lie: famine makes no distinction; when a city is taken, captivity desolates all; when the rocks tear the ship, ``the shipwreck is common without exception to all that sail in her.'' And then the twist that is not a consolation prize: ``if the Christian know and keep fast under what condition and what law he has believed, he will be aware that he must suffer \emph{more} than others in the world, since he must struggle more with the attacks of the devil.''\endnote{\emph{De mortalitate} 8--9, ANF V, read as above. The emphasis on ``more'' is this essay's; Cyprian's sentence carries it without typographic help. He supports the point from Ecclesiasticus 2:1--5 (``prepare your soul for temptation\ldots{} acceptable men in the furnace of humiliation'').}

That is the shape of the Christian answer to the two tenses, stated before anyone had a technical vocabulary for it. The facts of death are not altered. Nothing has been subtracted from the epidemic. What has changed is what the facts are \emph{for}.

\subsection{Why the baptized still die}

Augustine puts the same question in its sharpest analytical form and answers it with an argument that has no sentiment in it at all. ``If, moreover, any one is solicitous about this point, how, if death be the very punishment of sin, they whose guilt is cancelled by grace do yet suffer death''---his answer is that the parting of soul and body was deliberately left in place, ``for this reason, that if the immortality of the body followed immediately upon the sacrament of regeneration, faith itself would be thereby enervated. For faith is then only faith when it waits in hope for what is not yet seen in substance.''\endnote{Augustine, \emph{De civitate Dei} XIII.4, Dods 1871, chapter heading ``Why death, the punishment of sin, is not withheld from those who by the grace of regeneration are absolved from sin''; verified 2026-07-26 in the source library's tracked transcription of the 1871 printing, volume~1. Augustine refers the reader to his earlier work on the baptism of infants for the fuller treatment.}

He then presses it into the place where the argument earns its keep. If baptism removed death, the martyrs would have had nothing to give: ``there could not even have been any conflict, if, after the laver of regeneration, saints could not suffer bodily death. Who would not, then, in company with the infants presented for baptism, run to the grace of Christ, that so he might not be dismissed from the body?'' And so the inversion, which is one of the great sentences of Latin Christianity: ``Then it was proclaimed to man, `If thou sinnest, thou shalt die;' now it is said to the martyr, `Die, that thou sin not.'\,''\endnote{\emph{De civitate Dei} XIII.4, Dods 1871; verified as above. The chapter continues: ``That which was formerly set as an object of terror, that men might not sin, is now to be undergone if we would not sin\ldots{} For then death was incurred by sinning, now righteousness is fulfilled by dying.''}

The threat has been turned around and pointed the other way. The same instrument that enforced the prohibition is now the means of keeping it.

But Augustine will not let the inversion become an aesthetic. The very next chapter closes off the misreading with a symmetry: ``as the law is not an evil when it increases the lust of those who sin, so neither is death a good thing when it increases the glory of those who suffer it\ldots{} the law is indeed good, because it is prohibition of sin, and death is evil because it is the wages of sin; but as wicked men make an evil use not only of evil, but also of good things, so the righteous make a good use not only of good, but also of evil things.''\endnote{\emph{De civitate Dei} XIII.5, Dods 1871, chapter heading ``As the wicked make an ill use of the law, which is good, so the good make a good use of death, which is an ill''; verified as above. The chapter opens from 1 Corinthians 15:56 (``The sting of death is sin, and the strength of sin is the law'') and Romans 7:12--13.} And at the end of the earlier chapter, the statement that governs everything else in this section: ``Not that death, which was before an evil, has become something good, but only that God has granted to faith this grace, that death, which is the admitted opposite to life, should become the instrument by which life is reached.''\endnote{\emph{De civitate Dei} XIII.4 (closing), Dods 1871; verified as above. Compare XIII.7: ``But not on this account should we look upon death as a good thing, for it is diverted to such useful purposes, not by any virtue of its own, but by the divine interference.''}

Death has not been rehabilitated. It has been conscripted.

\subsection{What the Catechism keeps and what it does not}

The Catechism's treatment of Christian death is unusually warm, and it is worth checking whether the warmth costs anything. ``Because of Christ, Christian death has a positive meaning\ldots{} What is essentially new about Christian death is this: through Baptism, the Christian has already `died with Christ' sacramentally, in order to live a new life; and if we die in Christ's grace, physical death completes this `dying with Christ' and so completes our incorporation into him in his redeeming act.''\endnote{CCC 1010, Holy See English web text, read 2026-07-26, quoting Philippians 1:21 and 2 Timothy 2:11 and appending Ignatius of Antioch, \emph{Ad Romanos} 6, 1--2 (``It is better for me to die in Christ Jesus than to reign over the ends of the earth''). CCC 1009 supplies the christological ground: ``Jesus, the Son of God, also himself suffered the death that is part of the human condition. Yet, despite his anguish as he faced death, he accepted it in an act of complete and free submission to his Father's will. the obedience of Jesus has transformed the curse of death into a blessing.''} And CCC 1011, quoting Paul: ``the Christian can experience a desire for death like St. Paul's: `My desire is to depart and be with Christ.'\,''

Two things are worth noticing about the register. First, the Catechism attributes the transformation to Christ's obedience, not to death; the phrase is ``the obedience of Jesus has transformed the curse of death into a blessing,'' which locates the change exactly where Augustine located it---in the divine act, not in a property death has acquired. Second, the Catechism does not withdraw the harder statements. It is the same document that says death is a consequence of sin, contrary to the plan of the Creator, and the last enemy left to be conquered. The warmth is not purchased by suppressing the enmity.

Paul's own sentence, read at its own locus, has the same double structure. ``For to me, to live is Christ: and to die is gain. And if to live in the flesh: this is to me the fruit of labour. And what I shall choose I know not. But I am straitened between two: having a desire to be dissolved and to be with Christ, a thing by far the better. But to abide still in the flesh is needful for you.''\endnote{Philippians 1:21--24, Douay--Rheims, read 2026-07-26 in the tracked Challoner verse text. The Douay's ``to be dissolved'' renders \latin{dissolvi}; the verb is the same family as the ``dissolution'' of the earthly house at 2 Corinthians 5:1. The passage is quoted here for its structure---a preference stated and then declined---and the essay does not enter the exegetical question of what Paul expects to happen to him immediately upon dying, which belongs to the next section.} That is not a man in love with death. That is a man who has weighed two goods, found the further one better, and then set it down because he is needed. The desire is real and the desire is not decisive, and the pastoral value of the passage lies precisely in its refusal to resolve.

\subsection{Two tenses, two questions}

There is a tidy relation between the two tenses of the New Testament and the two questions of the previous section, and it is worth stating explicitly, as this essay's own reading rather than as anyone's doctrine.

Death was analyzed above under two heads: as the dissolution of a composite (a fact about a nature), and as the withdrawal of a preserving gift (a fact about a relationship). Christ's work bears on those two heads at different times. What is finished---what 2 Timothy can put in the past tense---is the second: for the man in Christ, death is no longer the enforcement of a sentence, because the sentence has been discharged, and this is true of him now, before he dies and whatever his prognosis. What is unfinished---what 1 Corinthians puts at the end of a queue---is the first: composites still come apart, and will until the resurrection makes the point moot by giving back what dissolution took.

So the celebrated ``already and not yet'' of Christian eschatology is not, at least here, a piece of vagueness or a way of postponing an embarrassment. It is the exact shape produced by the distinction the previous section defended. The penal meaning of death is cancelled in principle and the natural fact of it is not, because they were never the same thing, and there is no reason a single act should have to dispose of both at once.\endnote{This paragraph and the one preceding are project synthesis: an editorial reading connecting the two tenses of the New Testament witness to the natural/penal distinction defended in the previous section. It is argued from the texts cited in both sections and is attributed to no source. Its weakness is that it uses a scholastic distinction as a key to an apostolic text, and a reader who rejects the distinction has no reason to accept the reading. The claim is offered as an intelligible mapping, not as exegesis of Paul's intention.}

Which leaves standing the question the queue implies. If composites still come apart, and if the repair is at the end of the queue rather than at the moment of dissolution, then something is true of the interval---and the Church has said more about that interval, and defined more of it, than most Catholics realize.
